\documentclass[11pt]{article}
\usepackage[utf8]{inputenc}
\usepackage{amsmath}
\usepackage{amssymb}

% simple isotope helper (avoids the isotope package dependency)
\newcommand{\iso}[2]{\textsuperscript{#1}#2}

\begin{document}

\begin{center}
  {\large\textbf{Statistical thinning of pooled shards:}}\\[2pt]
  {\large\textbf{dose, acquisition start, washout, and per-species sampling}}
\end{center}
\medskip

\noindent
The range precision $\sigma_R$ --- the run-to-run spread of the reconstructed
endpoint $R_{50}$ at a given dose --- is a statistical quantity: it must be
estimated from an ensemble of acquisitions. Generating that ensemble by direct
simulation is expensive (each member is a full Monte Carlo plus reconstruction),
and the productions supply only a handful of independent realizations
(``shards''). Thinning is the technique that turns a small pool of shards into a
large ensemble of realizations, and, through the same operation, applies the
acquisition-start cut, the isotope washout, and the per-species selection. This
note sets out the method, why it is valid, what it buys, and where it fails.

\subsection*{The operation}

Let there be $n$ statistically independent shards, each a complete acquisition at
the reference dose (here $1$~Gy), produced with independent random seeds. Pool
them into one master list of $M$ recorded coincidences --- physically $n$~Gy of
events. A \emph{thinned realization} keeps each event $e$ of the pool
independently with a probability $q_e$: draw $u_e\sim\mathcal{U}[0,1)$ and keep
$e$ iff $u_e<q_e$. The kept subset is reconstructed exactly as a real
acquisition --- depth profile, erfc edge fit --- yielding one endpoint $R_{50}$.
Repeating with independent draws $N$ times gives $N$ values $\{R_{50}^{(1)},\dots,
R_{50}^{(N)}\}$, and
%
\begin{equation}
  \sigma_R \;=\; \mathrm{std}\bigl(R_{50}^{(1)},\dots,R_{50}^{(N)}\bigr),
  \qquad
  \frac{\delta\sigma_R}{\sigma_R}\;\simeq\;\frac{1}{\sqrt{2(N-1)}}.
  \label{eq:sigmaR}
\end{equation}
%
Everything specific to a study lives in the keep-probability $q_e$.

\subsection*{The keep-probability: dose, acquisition start, washout}

For a realization of dose $D$, acquisition start $t_s$, with washout, the
keep-probability factorises into independent effects,
%
\begin{equation}
  q_e \;=\; \underbrace{\mathbf{1}\!\left[t_{\mathrm{d}}(e)\ge t_s\right]}
              _{\text{acquisition start}}
        \;\times\; \underbrace{p_{\mathrm{dose}}}_{\text{dose}}
        \;\times\; \underbrace{w\bigl(z_0(e),t_{\mathrm{d}}(e)\bigr)}
              _{\text{washout}},
  \qquad
  p_{\mathrm{dose}} = \frac{D}{n},
  \label{eq:keep}
\end{equation}
%
where $z_0(e)$ is the event's origin depth and $t_{\mathrm{d}}(e)$ its decay
time, both recorded per coincidence.

\paragraph{Dose.}
Keeping a fraction $p_{\mathrm{dose}}=D/n$ of an $n$-Gy pool produces $D$~Gy of
events. This is the statistical sub-sampling that dials the dose and lets
$\sigma_R(D)$ be mapped without new simulation; scaling a measured $\sigma_R(D)$
to $1$~Gy uses $\sigma_R\propto 1/\sqrt{D}$.

\paragraph{Acquisition start.}
The hard cut $t_{\mathrm{d}}\ge t_s$ drops decays occurring before the scan opens;
a later start keeps fewer, later-decaying events.

\paragraph{Washout.}
$w(z_0,t_{\mathrm{d}})=\sum_i P(i\,|\,z_0,t_{\mathrm{d}})\,g_i$ is the
isotope-marginalised survival: $g_i$ is the per-species washout survival factor
and $P(i\,|\,z_0,t_{\mathrm{d}})$ the posterior that a recorded decay of depth
$z_0$ and decay time $t_{\mathrm{d}}$ belongs to species $i$ (both derived in the
companion washout note). Keeping each event with probability $w$ applies washout
as a per-event loss. Where the emitting species is known (as it is in the
production Monte Carlo), the marginalisation is unnecessary and $w$ is replaced
directly by the event's $g_i$.

Two variants replace the washout factor by a different third factor:
%
\begin{itemize}
  \item \emph{Nominal} (no washout): $w\equiv 1$.
  \item \emph{Per-species}: $w \to P(i\,|\,z_0,t_{\mathrm{d}})$ for a single
        species $i$, selecting an ``$i$-like'' sub-sample at its natural
        abundance --- used to measure the per-species precision
        $\sigma_R^{(i)}$.
\end{itemize}
%
A single machine --- a product of independent Bernoulli keep-factors on the
pooled list --- thus produces the dose sweep, the start-time axis, the washout
study, and the per-species study.

\subsection*{Why the sub-samples are valid, independent acquisitions}

Two facts, both consequences of the thinning theorem for Poisson processes,
underpin the method.

\paragraph{Thinned Poisson is Poisson.}
Radioactive production is a Poisson process, and independent Bernoulli thinning of
a Poisson process is again Poisson, with intensity scaled by the keep-probability.
In the large-pool limit a thinned realization is therefore a genuine lower-dose
(and/or washed) Poisson acquisition, not an approximation. This is why a
\emph{constant} per-species keep $g_i$ reproduces not only the mean washed profile
but its correct counting noise, with no per-event age required.

\paragraph{Realizations are mutually independent given the pool.}
Realization $A$ is a function only of its selection variables $\{u_e^{A}\}$ and
$B$ only of $\{u_e^{B}\}$; these are independent draws, and the pool positions
$\{z_e\}$ are fixed. Hence $R_{50}^{A}$ and $R_{50}^{B}$ are independent ---
sharing events does not correlate them, because the shared quantities are not
random within the estimate. The $N$ endpoints are therefore $N$ independent
samples of the $R_{50}$ distribution, and Eq.~(\ref{eq:sigmaR}) carries genuine
$N$-sample precision: $\pm10\%$ at $N=50$, $\pm5\%$ at $N=200$.

\subsection*{What thinning buys against direct shard statistics}

The direct alternative --- ``from-shards'' --- reconstructs each of the $n$ shards
once and takes $\sigma_R=\mathrm{std}(R_{50}^{(1)},\dots,R_{50}^{(n)})$ with error
$1/\sqrt{2(n-1)}$. For $n=10$ this is $\pm24\%$. It is statistically wasteful ---
it collapses each shard's millions of events into a single $R_{50}$ and estimates
$\sigma_R$ from $n$ numbers --- and fragile: at $n=10$ one atypical shard swings
the estimate (an $n=10$ artifact of exactly this kind produced a spurious
``washout is free'' reading early in this work). Thinning instead uses all the
pooled events to characterise the counting fluctuation, drawing $N\gg n$
independent realizations of it, and tightens $\sigma_R$ to $\pm10\%$ ($N=50$) or
$\pm5\%$ ($N=200$) --- enough to resolve configurations whose $\sigma_R$ differ at
the $10$--$20\%$ level, which $\pm24\%$ cannot.

The gain is real, not manufactured, because $\sigma_R$ is dominated by counting
noise, which the large pool determines precisely, and because the shards are
i.i.d.\ acquisitions of one fixed truth --- there is no shard-to-shard systematic
beyond counting. That last point is the single assumption the direct estimate
would protect against and thinning would not: were there a seed-dependent bias,
from-shards would see it in the spread of the $n$ endpoints while the thinned
realizations, sharing one pool, would not. Here there is none, so the tighter
estimate is honest.

\subsection*{Assumptions and regime of validity}

\begin{itemize}
  \item \emph{Poisson production; clearance independent per nucleus; spatially
        uniform survival} (needed for a position-independent $g_i$ and $w$). A
        non-uniform perfusion field breaks the last assumption and is the one
        route to a genuine, non-calibratable bias.
  \item \emph{Pool much larger than each realization} ($p_{\mathrm{dose}}$
        small). Sub-sampling a \emph{fixed} set of $M$ points is binomial
        (variance $\propto p(1-p)$), not Poisson ($\propto p$), so the thinned
        $\sigma_R$ is low by $\sqrt{1-p}$: negligible at $p\lesssim0.1$
        ($\le5\%$), but severe as $p\to1$, where every realization approaches the
        whole pool and $\sigma_R\to0$.
  \item \emph{Realization above the fit-stability floor} ($\sim\!300$--$400$k
        events). Below it the erfc edge fit destabilises on the noisy profile and
        $\sigma_R$ inflates spuriously (upward), not as a true spread.
  \item \emph{Thinning sharpens the spread, not the mean.} The pool's own edge
        differs from the truth by a $\sqrt{n}$-level fluctuation; that is the
        floor on the mean $R_{50}$ (the calibration / bias), which thinning does
        not reduce. Only $\sigma_R$ benefits.
\end{itemize}

\subsection*{The shrinking-pool corner}

The two count-dependent effects above pull in opposite directions, and at the
corner of small detector, long delay, and heavy washout --- where the surviving
events (pool $\times\,p_{\mathrm{dose}}\times w\times$ the start cut) collapse ---
they cannot be satisfied together. Clearing the fit floor demands a larger
$p_{\mathrm{dose}}$; keeping the fixed-pool bias $\sqrt{1-p}$ small demands a
smaller one. On a small pool there is no setting that does both, and $\sigma_R$
there is genuinely ill-determined: too large from fit instability, or too small
from the fixed-pool correction. The only clean remedy is a larger pool --- more
shards --- which buys headroom on both sides at once. A single-shard geometry
($p_{\mathrm{dose}}=D$ with no pool to thin) therefore cannot be treated this way,
and count-starved corners should be flagged rather than quoted at the nominal
$\pm10\%$.

\end{document}
